\section{Degree h=u+3: universal late-D-LC now pays a new actual C boundary}
The registered universal h>u late-D-LC theorem now supplies the
last derivative premise in this new actual C entrance. The resulting
ALL-padding late-C threshold is u>=h-3, with tail capacity/count and
binary count unbounded. Original C/A2/UI remain OPEN.

Let h=u+3, r=len(V), W=h+r, w=u+r. The actual first binary-count
loss from the infinite-binary same-tail reference is at N0=W-1,
with derivative loss z^N0 E/(1-z-z^2), E=(1+z)^u prod[t].
The separate initial-ones loss is z^(N0+alpha-1)E.
The entire reference prefix through W-2 is therefore retained.
Since capacities>=3, E0=1,E1=w,E2=binom(w,2)+r.
At ranks W-1,W,W+1 the EXACT total derivative losses (a,b,c) are
\[
 (2,2w+1,2E2+w+2)\ (alpha1),
\]
\[
 (1,w+2,E2+2w+2)\ (alpha2),\quad
 (1,w+1,E2+w+3)\ (alpha3),
\]
and (1,w+1,E2+w+2) for alpha>=4. At every w>=6,r>=1 they satisfy
1<=a<=2, b>=2a, c>=3b. The E2 coefficient includes the extra r;
neither finite-count loss may be omitted.

Write x=reference D_W, y=reference D_(W+1). Existing homogeneous
reference identities give 2<=x<=y<phi*x<2x and
A=x^2+x*y-y^2>0 in this range. The numerator has leading B_W=1,
so the exact ordinary reference reads are
\[
 Fref_(W-2)=x-1,\quad Fref_(W-1)=y,\quad Fref_W=x+y.
\]
The previous x-1 is essential; the homogeneous identity starts one
rank later and must not erase that numerator unit.
The ACTUAL values are consequently
\[
 D_W=x-b,\ D_(W+1)=y-c,\quad
 F_(W-2)=x-1,\ F_(W-1)=y-a,\ F_W=x+y-a-b.
\]
If the first actual tail pair is not strictly growing, universal
tail UI+LC makes all later C automatic. Otherwise lambda0=(y-c)/(x-b)>1.
Its golden gap is
phi*(x-b)-(y-c)>0 because c>=3b and y<phi*x.

The TWO ordinary comparisons needed to extend the unchanged cumulative
prefix are the following exact polynomials:
\[
 P=(y-a)(x-b)-(x-1)(y-c)
   =x(c-a)-y(b-1)+ab-c,
\]
\[
 Q=(x+y-a-b)(x-b)-(y-a)(y-c)
   =A-x(a+2b)+y(a-b+c)+b^2+ab-ac.
\]
For ANY real x>=2, x<=y<=2x, 1<=a<=2, b>=2a, c>=3b,
the first polynomial is >=2. Indeed use y<=2x and x>=2,
then c>=3b, to get P>=(a-1)(b-2)+2>=2.
For the second use y>=x, x>=a, c>=3b:
Q>=A+b(b-2a)>=A>0.
These are unbounded scalar implications, not host interpolation.

All G ratios through W-2 remain reference ratios and exceed phi.
P and Q extend the G comparison with lambda0 to W-1 and W.
For any later growing current N use lambda=D_N/D_(N-1)<=lambda0;
universal late-D-LC gives every earlier tail D ratio through N >=lambda.
The ordinary seed Q gives F_W>=lambda F_(W-1). Ordinary and cumulative
integration therefore pay every later read and EVERY shift>=3.
This proves all actual late C at h=u+3, all padding.

At w<6 one has u<=4, already covered by registered all-binary-counts
through eight. No new finite entrance certificate is required. Together
with the old h<=u+2 theorem, the new threshold is u>=h-3.
The remaining actual late-C residual becomes u>=9,u+4<=h<=3u-2,
maxcap<=2u-1,W>32. Earlier actual/reference C, A2 and original UI remain OPEN.

\paragraph{Formal and independent verification scope.}
Six new scoped real-algebra reports in \texttt{TripleSubdegreeCEntrance.lean}
check all-padding loss growth, both exact ordinary-seed identities,
both complete parameter-cone payments and the strict golden gap.
The source identification, reference numerator unit, mode/boundary
split and ordinary/cumulative integration are written assemblies
of the registered source/reference/universal-D-LC theorems, not an
end-to-end Lean composition theorem. All finite tests are QA ONLY;
there is NO new finite-host entrance premise or exception certificate.
